% ========================================================================
% INSTITUTO FEDERAL DO PIAUÍ - IFPI CAMPUS FLORIANO
% TEMPLATE DE PROJETO DE PESQUISA (ABNT NBR 15287 / MANUAL IFPI 2024)
% ========================================================================
% Este arquivo é a matriz principal do seu Projeto de Pesquisa.
% A estrutura é modular e importa os arquivos das pastas config/ e estrutura/.
%
% Conforme a ABNT NBR 14724:2024 (item 4.2.2), NBR 6024 e o Manual do IFPI,
% o trabalho é organizado exclusivamente em SEÇÕES (primárias, secundárias,
% terciárias, quaternárias e quinárias) e não em capítulos.
% ========================================================================

% ========================================================================
% CONFIGURAÇÕES E METADADOS
% ========================================================================
\documentclass[
    12pt,        % Tamanho da fonte principal
    openany,     % Seções iniciam em qualquer página
    oneside,     % Impressão apenas frente
    a4paper,     % Formato A4
    english,     % Idioma adicional
    french,      % Idioma adicional
    spanish,     % Idioma adicional
    brazil,      % Idioma principal
    article      % Estruturação por seções (sem \chapter) conforme NBR 14724:2024
]{abntex2}

\input{config/config}   % configurações de formatação
% ========================================================================
% DADOS DO PROJETO DE PESQUISA
% Edite os valores abaixo para personalizar o template.
% ========================================================================

% Instituição e identificação do curso
\instituicao{Instituto Federal de Educação, Ciência e Tecnologia do Piauí}
\campus{Campus Floriano}
\curso{Tecnologia em Análise e Desenvolvimento de Sistemas}
\areaconcentracao{} % Opcional: defina a área de concentração ou deixe vazio

% Título e autoria
\titulo{DESENVOLVIMENTO DE UMA PLATAFORMA WEB PARA CATALOGAÇÃO E CONSULTA DO ACERVO DO HERBÁRIO DO IFPI – CAMPUS FLORIANO}
\subtitulo{}
\tituloestrangeiro{Development of a Web Platform for Cataloging and Accessing the Herbarium Collection at IFPI – Floriano Campus}
\autor{ANDRE FELIPE OLIVEIRA DUTRA}
\emailautor{autor@aluno.ifpi.edu.br}
\vinculoautor{IFPI -- Campus Floriano}

% Orientação
\orientador{Prof. Me. Ronaldo Pires Borges}
\emailorientador{ronaldo.borges@ifpi.edu.br}
\vinculoorientador{IFPI -- Campus Floriano} % Vínculo do orientador

% Coorientação (deixe vazio se não houver coorientador)
\coorientador{}
\emailcoorientador{}
\vinculocoorientador{}

% Informações gerais
\local{Floriano -- PI}
\data{2026}
\dataaprovacao{XX/XX/XXXX} % ESCREVA A DATA DA BANCA / APROVAÇÃO

% Descrição apresentada na folha de rosto
\preambulo{Projeto de pesquisa apresentado ao Instituto Federal de Educação, Ciência e Tecnologia do Piauí -- Campus Floriano, como requisito para aprovação na disciplina de Projeto de Pesquisa do curso de Tecnologia em Análise e Desenvolvimento de Sistemas.}

% Natureza do trabalho
% \naturezatrabalho{Projeto de pesquisa apresentado ao Instituto Federal de Educação, Ciência e Tecnologia do Piauí como requisito para aprovação...}

% Grau pretendido (opcional)
% \grau{Tecnólogo em Análise e Desenvolvimento de Sistemas} 

% ========================================================================
% INSTRUÇÕES RÁPIDAS
% - Substitua os textos pelos dados reais da sua pesquisa.
% - Campos opcionais sem uso podem permanecer vazios.
% ========================================================================
 % dados institucionais e do projeto

% ========================================================================
% INÍCIO DO DOCUMENTO
% ========================================================================

\begin{document}
% Seleciona o idioma do documento (conforme pacotes do babel)
\selectlanguage{brazil}
% Retira espaço extra obsoleto entre as frases
\frenchspacing

% ========================================================================
% ELEMENTOS PRÉ-TEXTUAIS
% ========================================================================
\input{estrutura/pre_textuais}

% ========================================================================
% ELEMENTOS TEXTUAIS
% ========================================================================

\section{INTRODUÇÃO}
\label{sec:introducao}

Os herbários reúnem exemplares botânicos preservados e informações associadas à sua coleta e identificação. Esses acervos apoiam atividades de ensino e pesquisa, pois permitem consultar registros de espécies e estudar a diversidade vegetal. No Instituto Federal do Piauí (IFPI), Campus Floriano, o herbário também constitui um espaço de aprendizagem prática para estudantes de Ciências Biológicas, conforme descrito por \citeonline{lopes2023herbario}.

A organização das informações de um herbário é necessária para que os exemplares possam ser localizados e consultados. \citeonline{teixeira2016herbario} discute a importância do registro e da informatização das exsicatas para o controle do acervo e o acesso às informações. Estudos sobre a digitalização de coleções botânicas também apontam que a disponibilização de dados e imagens pode ampliar as possibilidades de consulta e reduzir a necessidade de manuseio dos exemplares físicos \cite{thompson2023digitisation}.

Nesse contexto, este trabalho propõe o desenvolvimento de uma plataforma web para apoiar a catalogação e a consulta do acervo do Herbário do IFPI – Campus Floriano. A proposta concentra-se em um conjunto delimitado de funcionalidades, como o cadastro e a consulta de registros de exsicatas, com informações de identificação e coleta e, quando disponíveis, imagens dos exemplares. A definição dos campos e das funcionalidades deverá considerar as necessidades dos responsáveis pelo herbário.

Espera-se que a plataforma ofereça uma forma organizada de registrar e acessar as informações do acervo e sirva como base para futuras ampliações. Assim, o trabalho articula o desenvolvimento de software a uma necessidade concreta do campus, com foco em uma solução de alcance compatível com o projeto de conclusão de curso.


\subsection{PROBLEMA DE PESQUISA}
\label{subsec:problema}

O problema de pesquisa constitui o elemento central que orienta todo o desenvolvimento do estudo. Conforme \citeonline{gil2019}, um problema bem formulado deve ser claro, preciso e operacional, permitindo a definição de objetivos específicos e a escolha de métodos adequados para sua investigação.

No contexto desta pesquisa, observa-se que as informações sobre as exsicatas do Herbário do IFPI -- Campus Floriano estão associadas aos próprios exemplares e a um livro de registro \cite{lopes2023herbario}, o que torna os dados dependentes da conservação do material físico e restringe sua consulta ao acesso presencial. Esta situação tem gerado a perda de exemplares por deterioração, como a de cerca de 300 exsicatas descartadas quando as atividades foram retomadas em 2018 \cite{lopes2023herbario}, evidenciando a necessidade de estudos que possam contribuir para sua compreensão e possível solução.

Diante do exposto, formula-se o seguinte problema de pesquisa:

\textbf{Como uma plataforma web pode armazenar de forma organizada as informações das exsicatas do Herbário do IFPI -- Campus Floriano e disponibilizar o catálogo do acervo para consulta por visitantes?}

Este questionamento surge da observação de que a informatização do acervo é apontada como meio de controlar as coleções e ampliar o acesso às informações \cite{teixeira2016herbario}, conforme evidenciado por \citeonline{thompson2023digitisation} em seus estudos sobre o fluxo de digitalização em um herbário universitário.


\subsection{OBJETIVOS}
\label{subsec:objetivos}

Os objetivos definem o que se pretende alcançar com a realização da pesquisa. Devem ser formulados de forma clara e precisa, utilizando verbos no infinitivo que expressem ações concretas e mensuráveis.

\subsubsection{Objetivo Geral}
\label{subsubsec:objetivo-geral}

Desenvolver uma plataforma web para o registro e a organização das exsicatas do Herbário do IFPI -- Campus Floriano e para a consulta pública do catálogo do acervo por visitantes.

\subsubsection{Objetivos Específicos}
\label{subsubsec:objetivos-especificos}

\begin{enumerate}
    \item Compreender o processo de coleta, identificação, armazenamento e consultas das exicatas do acervo; 
    \item Modelar o banco de dados a partir das informações levantadas e das práticas de curadoria de herbários;
    \item Implementar a área de gerenciamento do acervo, com cadastro, edição e associação de imagens às exsicatas;
    \item Implementar a área pública de consulta ao catálogo, com busca e visualização dos registros;
    \item Validar a plataforma com os responsáveis pelo herbário e com visitantes.
\end{enumerate}


\subsection{JUSTIFICATIVA}
\label{subsec:justificativa}

A justificativa deve demonstrar a relevância e a importância da pesquisa, apresentando argumentos que evidenciem sua contribuição para o avanço do conhecimento científico e/ou para a solução de problemas práticos.

Esta pesquisa justifica-se por diversos aspectos:

\textbf{Relevância Científica:} O tema abordado contribui para o avanço do conhecimento na área de [ÁREA DE CONHECIMENTO], especialmente no que se refere a [ASPECTO ESPECÍFICO]. Conforme destacado por \citeonline{inep2022}, existe uma lacuna na literatura sobre [LACUNA IDENTIFICADA], que esta pesquisa pretende contribuir para preencher.

\textbf{Relevância Social:} Os resultados desta pesquisa poderão beneficiar [PÚBLICO-ALVO OU SOCIEDADE EM GERAL], uma vez que [EXPLICAR COMO OS RESULTADOS PODEM SER APLICADOS]. Segundo dados de \cite{instituicao2023}, [APRESENTAR DADOS QUE EVIDENCIEM A RELEVÂNCIA SOCIAL].

\textbf{Relevância Acadêmica:} Este estudo contribuirá para a formação acadêmica do pesquisador e para o desenvolvimento de pesquisas futuras na área, fornecendo subsídios teóricos e metodológicos para estudos similares.

\textbf{Viabilidade:} A pesquisa é viável considerando os recursos disponíveis, o tempo previsto para sua execução e o acesso às fontes de informação necessárias.


\section{REFERENCIAL TEÓRICO}
\label{sec:referencial}

O referencial teórico constitui a base conceitual da pesquisa, apresentando as teorias, conceitos e estudos que fundamentam a investigação. Conforme exposto por \citeonline{IFPI_Manual_2024}, esta seção destina-se a detalhar as leituras e embasamentos teóricos que sustentam a proposta, explicando, classificando e demonstrando o estado da arte do objeto de estudo.

\subsection{HERBÁRIO: CONCEITO E IMPORTÂNCIA}
\label{subsec:conceito-herbario}

Um herbário é uma coleção científica de plantas ou fungos secos, fixados em cartolina e preservados para fins científicos. Cada exemplar acompanha uma etiqueta com dados como o local de coleta, características da planta e seus usos, e todo o material preservado dessa forma recebe o nome de exsicata \cite{lopes2022herbarios}. Em sentido mais amplo, os herbários são coleções biológicas de plantas, fungos e algas pluricelulares com finalidades científicas e educacionais, essenciais para o estudo da taxonomia botânica, da distribuição geográfica e da nomenclatura dos grupos botânicos \cite{teixeira2016herbario}.

\citeonline{peixoto2003colecoes} classificam os herbários como coleções preservadas, em oposição às coleções vivas, como os jardins botânicos, e afirmam que são ferramentas imprescindíveis para o trabalho dos taxonomistas e apoio para muitas outras áreas do conhecimento. Segundo os autores, o herbário fornece o material-testemunho (\textit{voucher}) de um organismo, constitui a base de dados sobre a distribuição geográfica e a diversidade das plantas e guarda a memória dos conceitos morfológicos e taxonômicos e das mudanças pelas quais passaram. À época do estudo, havia 114 herbários ativos no Brasil, e o \textit{Index Herbariorum} listava 3.210 herbários no mundo, dos quais 73 brasileiros \cite{peixoto2003colecoes}.

A importância dessas coleções cresceu após a Convenção sobre a Diversidade Biológica, assinada em 1992 \cite{teixeira2016herbario,peixoto2003colecoes}. \citeonline{souza2022herbarios} apontam que os herbários abrigam amostras da biodiversidade vegetal de diferentes ecossistemas e apoiam pesquisas em áreas como filogenia, taxonomia e fenologia, além de formarem um sistema de informação sobre plantas locais ou exóticas, inclusive daquelas que já não são encontradas em determinada região. \citeonline{teixeira2016herbario} acrescenta que as coleções biológicas podem atender à demanda por informações para a avaliação de impactos ambientais, a definição de áreas de preservação, a proteção de espécies ameaçadas e o estabelecimento de políticas públicas.

Os herbários também têm função didática. \citeonline{lopes2022herbarios} revisaram teses e dissertações brasileiras defendidas entre 2011 e 2020 sobre o uso de herbários no aprendizado de Botânica e concluíram que o tema ainda é pouco estudado, mas que o herbário didático desperta a curiosidade científica dos estudantes e incentiva a conscientização ambiental. No contexto do IFPI Campus Floriano, \citeonline{lopes2023herbario} relata que o contato com as exsicatas permite ao discente identificar, por meio das etiquetas, nomes científicos, datas e informações sobre a coleta e a localização das plantas, o que desperta o interesse pela taxonomia e pela conservação das espécies.

O Herbário do IFPI Campus Floriano teve origem nas primeiras exsicatas produzidas em uma visita técnica realizada em 6 de novembro de 2003 e, em 2007, passou a contar com um livro de registro, que reunia 526 exsicatas \cite{lopes2023herbario}. Com a interrupção das atividades, as coordenadoras retomaram o trabalho em 2018 e constataram a perda de cerca de 300 exsicatas, descartadas por deterioração; na data do relato, o herbário possuía cerca de 200 exsicatas \cite{lopes2023herbario}. Esse histórico evidencia que a conservação do acervo e a preservação das informações a ele associadas são desafios concretos para a instituição.

\subsection{COLETA, HERBORIZAÇÃO E CATALOGAÇÃO DO ACERVO}
\label{subsec:coleta-catalogacao}

% TODO (André): pedir à equipe do Herbário os manuais/protocolos que eles seguem
% (ficha de coleta, livro de registro, regras de tombamento, modelo de etiqueta).
% Esta subseção foi escrita com base no guia do INSA e nos trabalhos de Teixeira (2016)
% e Lopes (2023); quando receber os manuais, complemente/ajuste os trechos abaixo
% e acrescente as referências no referencias.bib.

A obtenção de um exemplar de herbário segue um conjunto de procedimentos, que vai da coleta em campo à catalogação, e que assegura a qualidade da pesquisa e a credibilidade científica da coleção \cite{freitas2021coleta}. Esta subseção resume essas etapas com base em um guia prático de coleta de material botânico e em trabalhos sobre a manutenção de herbários didáticos; os procedimentos específicos adotados pela equipe do Herbário do IFPI Campus Floriano serão complementados com os manuais utilizados por ela.

\subsubsection{Coleta em campo}

A coleta de plantas e fungos no Brasil exige licença ou autorização prévia das instituições competentes, e o material biológico nacional utilizado em pesquisa deve ser registrado no Sistema Nacional de Gestão do Patrimônio Genético e do Conhecimento Tradicional Associado (SisGen) \cite{freitas2021coleta}. Sempre que possível, a amostra deve conter estruturas vegetativas e reprodutivas, que são as mais úteis para a identificação botânica, e recomenda-se obter duplicatas de duas a cinco amostras da mesma planta \cite{freitas2021coleta}. O coletor deve registrar no próprio local o máximo de informações sobre as amostras, especialmente as que se perdem na herborização, como odores e cores, e as que não são evidentes no material, como o habitat e a relação com outros organismos \cite{freitas2021coleta}.

\subsubsection{Herborização}

A herborização compreende a prensagem, a secagem e a montagem das exsicatas \cite{freitas2021coleta}. Na prensagem, as amostras são acomodadas entre folhas de papel ou jornal e papelão, com o número de coleta identificado, e comprimidas entre duas grades de madeira amarradas; em seguida, as prensas são levadas a uma estufa, onde o material é desidratado, e devem ser reapertadas periodicamente, porque as amostras perdem água e geram folga na prensa \cite{freitas2021coleta}. Na montagem, a amostra seca é fixada, por costura ou colagem, em cartolina livre de ácido de $30 \times 42$ cm e recebe uma etiqueta de $10 \times 12$ cm com a sigla do herbário, a identificação botânica, a data e os dados de origem, além do carimbo com o número de tombo do herbário \cite{freitas2021coleta}. Amostras adicionais com o mesmo número de coleta são chamadas duplicatas e servem a intercâmbios entre herbários; quando não há duplicatas, a exsicata é considerada unicata \cite{freitas2021coleta}.

\subsubsection{Catalogação e manutenção}

Depois de herborizadas e incorporadas à coleção, as plantas passam a ser denominadas exsicatas e recebem um número de registro (tombo), identificado por um carimbo da instituição; cada exsicata possui uma ficha com o local da coleta, o coletor, o habitat, usos populares ou medicinais e o número de coleta \cite{teixeira2016herbario}. A manutenção da coleção inclui a verificação da presença de fungos e a degradação das embalagens, a desinsetização dos exemplares e dos armários e a revisão das fichas de identificação para corrigir erros de identificação e de agrupamento em famílias \cite{teixeira2016herbario}. Novas exsicatas podem ser identificadas por comparação com o material já identificado, por chaves dicotômicas e por literatura específica \cite{teixeira2016herbario}.

No Herbário do IFPI Campus Floriano, as coordenadoras adquiriram por conta própria um curso de introdução à taxonomia de campo para treinar os discentes nas atividades de herborização, e o registro dos exemplares era mantido em um livro de registro criado em 2007 \cite{lopes2023herbario}. O \autoref{qua:etapas} sintetiza as etapas descritas e as informações geradas em cada uma delas, que são justamente os dados que uma plataforma de catalogação precisa representar.

\begin{quadro}[htbp]
    \centering
    \caption{Etapas do tratamento do material botânico e informações geradas}
    \label{qua:etapas}
    \begin{tabularx}{\textwidth}{|p{3cm}|X|X|}
        \hline
        \textbf{Etapa} & \textbf{O que é feito} & \textbf{Informações geradas} \\
        \hline
        Coleta & Obtenção de amostras com estruturas vegetativas e reprodutivas, com registro em campo & Número de coleta, local, data, coletor, habitat, cor e odor \\ \hline
        Herborização & Prensagem, secagem em estufa e montagem da exsicata & Exsicata montada com etiqueta de identificação \\ \hline
        Identificação & Comparação com material identificado, chaves dicotômicas e literatura & Nome científico e família \\ \hline
        Tombamento & Registro no acervo com carimbo da instituição & Número de tombo e registro no livro \\ \hline
        Manutenção & Inspeção, desinsetização e revisão das fichas & Correções de identificação e de agrupamento \\ \hline
    \end{tabularx}
    \fonte{Elaborado pelo autor, com base em \citeonline{freitas2021coleta}, \citeonline{teixeira2016herbario} e \citeonline{lopes2023herbario}.}
\end{quadro}

\subsection{ACESSO AO ACERVO E INFORMATIZAÇÃO}
\label{subsec:acesso-informatizacao}

Tradicionalmente, o acesso às informações de uma exsicata exige o manuseio do exemplar herborizado, o que torna o trabalho do pesquisador demorado e pode danificar o material \cite{teixeira2016herbario}. Para estudantes e pesquisadores, isso significa consultar o acervo presencialmente e depender do estado de conservação das exsicatas. \citeonline{souza2022herbarios} destacam que o herbário também permite a alunos do ensino fundamental, médio e da graduação conhecer o funcionamento da coleção e a importância das coleções científicas, e que a digitalização das exsicatas é uma das alternativas dos herbários para alcançar maior público.

A informatização dos herbários responde a esse problema. \citeonline{peixoto2003colecoes} já assinalavam que a informatização agiliza o processo de geração do conhecimento e que o georreferenciamento das amostras constitui informação relevante, mas que o processo ainda era incipiente no Brasil: apenas 52\% dos herbários tinham mais da metade do acervo informatizado, e 11\% sequer haviam começado. Para \citeonline{souza2022herbarios}, a implantação de um banco de dados agiliza o intercâmbio de informações entre regiões e reduz o envio de material físico, que em alguns casos danificava as amostras; além disso, a disponibilização on-line permite analisar os dados sem viagens exaustivas e contorna a burocracia de acesso às coleções. \citeonline{teixeira2016herbario} indica que o registro (tombo) e a informatização dos dados geram maior controle sobre o acervo e viabilizam a digitalização, em que cada exemplar é fotografado. \citeonline{thompson2023digitisation} também examinam o fluxo de digitalização em um herbário universitário.

\citeonline{souza2022herbarios}, citando Hopkins (2005), descrevem três usos da virtualização de um herbário: o uso intra-herbário, de consulta a imagens e dados do próprio herbário, que facilita a rotina de curadoria e reduz o risco de danos em consultas físicas; o uso inter-herbário, de consulta a dados de outros herbários para melhorar a curadoria local; e o uso extra-herbário, de consulta a vários herbários para fins que vão além da curadoria, como listas regionais de espécies e informações sobre coletores. A plataforma proposta nesta pesquisa se enquadra, em primeiro lugar, no uso intra-herbário, pois organiza e disponibiliza o acervo do próprio Herbário do IFPI Campus Floriano.

Para que os dados possam ser compartilhados, é útil adotar padrões. O \textit{Darwin Core} é um padrão do \textit{Biodiversity Information Standards} (TDWG) que oferece um conjunto de termos para descrever e compartilhar dados sobre a ocorrência de organismos e sobre os espécimes de coleções biológicas \cite{tdwg2009darwincore}. Adotá-lo como referência para a modelagem dos campos da plataforma facilita uma eventual integração futura com redes de herbários.
% VERIFICAR: confirmar no site do TDWG (https://dwc.tdwg.org) a descrição do padrão
% e, se for citar termos específicos (ex.: recordedBy, catalogNumber), conferir os nomes.

\subsection{PROCESSO DE DESENVOLVIMENTO DE SOFTWARE}
\label{subsec:processo-software}

Um processo de software é um conjunto de atividades que levam à produção de um sistema; em geral, essas atividades abrangem a especificação (levantamento e análise de requisitos), o projeto e a implementação, a validação (testes) e a evolução do software \cite{sommerville2011}. As organizações escolhem e combinam essas atividades de modos diferentes, o que dá origem a modelos de processo mais sequenciais ou mais iterativos.
% VERIFICAR: conferir com o livro (edição, capítulo 2 - Processos de software) antes de entregar.

Entre as abordagens iterativas, o \textit{Scrum} é um arcabouço leve que ajuda pessoas, equipes e organizações a gerar valor por meio de soluções adaptativas para problemas complexos \cite{schwaber2020scrum}. Ele se baseia no empirismo, ou seja, na tomada de decisão a partir do que se observa, e em três pilares: transparência, inspeção e adaptação \cite{schwaber2020scrum}. O trabalho é organizado em \textit{Sprints}, ciclos de no máximo um mês em que se produz um incremento utilizável do produto, e cada \textit{Sprint} contém eventos como o planejamento (\textit{Sprint Planning}), a reunião diária (\textit{Daily Scrum}), a revisão (\textit{Sprint Review}) e a retrospectiva (\textit{Sprint Retrospective}) \cite{schwaber2020scrum}. Os requisitos ficam em uma lista ordenada, o \textit{Product Backlog}, de onde se seleciona o que será feito em cada \textit{Sprint} (\textit{Sprint Backlog}) \cite{schwaber2020scrum}.

Esse tipo de processo é adequado a este projeto porque os requisitos do herbário, como os campos da ficha da exsicata e a forma de consulta pelos usuários, não estão totalmente definidos no início e tendem a ser refinados à medida que a equipe do herbário avalia os primeiros incrementos. O próprio guia reconhece que a equipe do Scrum é pequena, tipicamente com dez pessoas ou menos \cite{schwaber2020scrum}, o que permite adaptar o arcabouço a um projeto acadêmico individual, com o orientador e a equipe do herbário atuando como fontes dos requisitos e da avaliação dos incrementos.
% VERIFICAR/ALINHAR com o orientador: (1) se o Scrum será de fato o processo adotado
% (a seção METODOLOGIA ainda não o menciona); (2) os papéis acima; (3) a duração das Sprints.

\subsection{TECNOLOGIAS ESCOLHIDAS E ALTERNATIVAS}
\label{subsec:tecnologias}

A plataforma será construída como uma aplicação web, com \textit{back-end} em Node.js e TypeScript, banco de dados PostgreSQL e \textit{front-end} em React. Esta subseção descreve cada tecnologia e justifica a escolha em relação a alternativas.

O Node.js é um ambiente de execução JavaScript assíncrono, orientado a eventos, projetado para construir aplicações de rede escaláveis \cite{nodejs2026about}. O TypeScript é uma linguagem que adiciona tipos estáticos ao JavaScript e é convertida em JavaScript comum, o que permite detectar erros de tipo antes da execução \cite{typescript2026}. O React é uma biblioteca para a construção de interfaces de usuário a partir de componentes reutilizáveis \cite{react2026learn}. O PostgreSQL é um sistema gerenciador de banco de dados relacional de código aberto \cite{postgresql2026about}, e a extensão PostGIS acrescenta a ele tipos e funções para dados geográficos \cite{postgis2026}.

Os motivos da escolha são os seguintes. Primeiro, o uso de JavaScript/TypeScript no \textit{back-end} e no \textit{front-end} permite trabalhar com uma única linguagem em toda a aplicação, o que reduz a curva de aprendizado em um projeto desenvolvido por uma única pessoa e viabiliza o compartilhamento de tipos entre as camadas. Segundo, os dados de um herbário são fortemente relacionados (por exemplo, exsicata, espécie, gênero, família, coletor e local de coleta) e exigem unicidade do número de tombo, características atendidas de modo natural por um banco de dados relacional. Terceiro, o PostGIS permite armazenar e consultar as coordenadas de coleta, informação que \citeonline{peixoto2003colecoes} destacam como relevante nas coleções; essa combinação também é usada no sistema Jabot do Jardim Botânico do Rio de Janeiro, que emprega PostgreSQL e PostGIS para validar dados geográficos de coletas \cite{silva2019agrupamentos}. Quarto, todas as tecnologias são de código aberto e não exigem pagamento de licenças, condição importante para uma instituição pública de ensino.

O \autoref{qua:tecnologias} compara as tecnologias escolhidas com alternativas usuais para cada camada.

\begin{quadro}[htbp]
    \centering
    \caption{Tecnologias escolhidas e alternativas consideradas}
    \label{qua:tecnologias}
    \begin{tabularx}{\textwidth}{|p{2.2cm}|p{2.8cm}|p{3cm}|X|}
        \hline
        \textbf{Camada} & \textbf{Escolha} & \textbf{Alternativas} & \textbf{Justificativa} \\
        \hline
        \textit{Back-end} & Node.js com TypeScript & PHP (Laravel), Python (Django) & Mesma linguagem do \textit{front-end}; tipagem estática reduz erros em tempo de desenvolvimento \\ \hline
        Banco de dados & PostgreSQL com PostGIS & MySQL, MongoDB & Modelo relacional com integridade referencial; suporte a dados geográficos; uso em sistema semelhante (Jabot) \\ \hline
        \textit{Front-end} & React & Angular, Vue.js & Interface baseada em componentes reutilizáveis; mesma linguagem do \textit{back-end} \\ \hline
    \end{tabularx}
    \fonte{Elaborado pelo autor, com base em \citeonline{nodejs2026about}, \citeonline{typescript2026}, \citeonline{react2026learn}, \citeonline{postgresql2026about}, \citeonline{postgis2026} e \citeonline{silva2019agrupamentos}.}
\end{quadro}

% VERIFICAR: as alternativas do quadro (Laravel, Django, MySQL, MongoDB, Angular, Vue.js)
% foram citadas só como opções usuais, sem fonte consultada. Se quiser argumentar que uma
% delas é pior em algum ponto, busque referência (documentação oficial ou livro) para sustentar.

\subsection{ESTUDOS E SISTEMAS RELACIONADOS}
\label{subsec:estudos-relacionados}

Existem sistemas voltados à informatização e à disponibilização de acervos de herbários. O Herbário Virtual Reflora, resultado do Programa Reflora, iniciado em 2010, reúne imagens em alta resolução de exsicatas e dados associados repatriados de herbários estrangeiros, com o objetivo de reproduzir em ambiente virtual o que se pode fazer em um herbário real; foi lançado em 2013, atingiu um milhão de imagens em 2015 e pode ser consultado tanto pela comunidade botânica quanto pelo público em geral \cite{forzza2015reflora}. O Jabot, desenvolvido no Jardim Botânico do Rio de Janeiro entre 2003 e 2005 com \textit{software} livre (PHP e PostgreSQL), gerencia o herbário da instituição, que está totalmente informatizado e digitalizado, com cerca de 760 mil registros, e usa PostGIS para validação geográfica \cite{silva2019agrupamentos}. Outra solução é o Symbiota, plataforma exclusivamente web para a criação de comunidades de informação sobre biodiversidade baseadas em espécimes-testemunho (\textit{voucher}), que oferece amplos recursos de troca de dados \cite{gries2014symbiota}.
Há ainda iniciativas que integram vários acervos. O Instituto Nacional de Ciência e Tecnologia Herbário Virtual da Flora e dos Fungos (INCT-HVFF), aprovado em 2008 e em atividade desde 2009, tem como objetivos disponibilizar dados de herbários de forma livre, aberta e integrada, repatriar dados de coletas brasileiras depositadas no exterior e melhorar a qualidade dos acervos dos herbários do país, e disponibiliza o Herbário Virtual por meio da rede \textit{speciesLink} \cite{incthvff2026}. O Herbário Virtual Glaziou usa essa mesma rede para reunir informações e imagens das exsicatas coletadas por Auguste Glaziou no Brasil e depositadas no Museu Nacional do Rio de Janeiro e no Museu Nacional de História Natural de Paris \cite{glaziou2026}. Fora do Brasil, o Herbario Digital chileno oferece catálogo taxonômico, busca de imagens, seção de espécimes e mapa sobre a flora vascular do país \cite{herbariodigital2026}.
% VERIFICAR: as páginas do INCT-HVFF, Glaziou e Herbario Digital não informam o número
% de registros; se quiser números, abra os sites no navegador e confira (data de acesso 12 out. 2026).
% O artigo do Jabot em si (Rodriguésia) não teve autoria/ano confirmados, por isso não foi incluído.

O \autoref{qua:relacionados} compara esses sistemas com a plataforma proposta.

\begin{quadro}[htbp]
    \centering
    \caption{Comparação entre sistemas relacionados e a plataforma proposta}
    \label{qua:relacionados}
    \begin{tabularx}{\textwidth}{|p{2.6cm}|p{3.2cm}|X|X|}
        \hline
        \textbf{Sistema} & \textbf{Escopo} & \textbf{Características relatadas} & \textbf{Diferença em relação à proposta} \\
        \hline
        Herbário Virtual Reflora \cite{forzza2015reflora} & Imagens e dados de exsicatas de herbários nacionais e estrangeiros & Alta resolução; consulta pública; mais de um milhão de imagens em 2015 & Escala nacional e internacional; não gerencia o livro de tombo nem a rotina de curadoria do herbário do IFPI \\ \hline
        Jabot \cite{silva2019agrupamentos} & Gerenciamento de coleções do Jardim Botânico do Rio de Janeiro & PHP e PostgreSQL; PostGIS para validação geográfica; cerca de 760 mil registros & Sistema institucional de grande porte; a proposta é uma plataforma enxuta, voltada a um herbário didático \\ \hline
        Symbiota \cite{gries2014symbiota} & Plataforma web para comunidades de dados de biodiversidade & Baseada em espécimes-testemunho; recursos de troca de dados & Plataforma genérica e configurável; a proposta é desenvolvida sob medida para o acervo do campus \\ \hline
        INCT-HVFF e \textit{speciesLink} \cite{incthvff2026} & Rede nacional de dados de herbários brasileiros e de coletas depositadas no exterior & Dados livres, abertos e integrados; repatriamento de dados; Herbário Virtual via \textit{speciesLink} & Integra dezenas de instituições; a proposta atende a um único herbário e à sua rotina de cadastro \\ \hline
        Herbário Virtual Glaziou \cite{glaziou2026} & Exsicatas coletadas por Glaziou, de dois museus & Informações e imagens das exsicatas; usa a rede \textit{speciesLink} & Acervo histórico e temático; a proposta cobre todo o acervo do herbário do campus \\ \hline
        Herbario Digital (Chile) \cite{herbariodigital2026} & Flora vascular do Chile & Catálogo taxonômico, busca de imagens, espécimes e mapa & Foco em outro país; serve de referência de organização da consulta \\ \hline
    \end{tabularx}
    \fonte{Elaborado pelo autor, com base nas páginas e nos autores citados.}
\end{quadro}

Esses sistemas mostram que a informatização de herbários é uma prática consolidada e que o PostgreSQL, tecnologia escolhida para esta proposta, é usado em um sistema reconhecido da área, o Jabot. Eles também deixam uma lacuna que a presente proposta busca preencher: são soluções de grande porte ou de propósito geral, voltadas a instituições de pesquisa, enquanto o Herbário do IFPI Campus Floriano é um herbário didático de acervo pequeno, cujo registro ainda depende de um livro e dos próprios exemplares \cite{lopes2023herbario}. A proposta se diferencia por oferecer uma plataforma simples, mantida pela própria instituição, que concentra o cadastro e a consulta das exsicatas e que pode, no futuro, adotar padrões como o \textit{Darwin Core} para facilitar o intercâmbio de dados com redes maiores.
% VERIFICAR com o orientador: confirmar se "enxuta/sob medida" é o posicionamento que
% você quer defender e se o texto não promete funcionalidades além do escopo dos objetivos.


\section{METODOLOGIA}
\label{sec:metodologia}

A metodologia descreve detalhadamente os procedimentos que serão adotados para a realização da pesquisa. O texto deve manter um tratamento impessoal e objetivo em conformidade com o \citeonline{IFPI_Manual_2024}.

\subsection{CARACTERIZAÇÃO DA PESQUISA}
\label{subsec:caracterizacao}

Esta pesquisa pode ser caracterizada, quanto aos seus objetivos, como EXPLORATÓRIA, uma vez que busca proporcionar maior familiaridade com o problema. Segundo \citeonline{gil2019}, pesquisas dessa natureza têm como objetivo o aprimoramento de ideias ou a descoberta de intuições.

Quanto à abordagem, trata-se de uma pesquisa QUANTITATIVA, pois adota métodos estatísticos para mensurar os dados coletados. \citeonline{creswell2018} destaca que a abordagem mista ou quantitativa é adequada quando se pretende testar teorias objetivas examinando a relação entre variáveis.

Em relação aos procedimentos técnicos, será realizada uma PESQUISA BIBLIOGRÁFICA, que, conforme \citeonline{marconi2017}, caracteriza-se pelo levantamento de contribuições teóricas já publicadas sobre o tema.

\begin{figure}[htbp]
    \centering
    \caption{Esquema metodológico da pesquisa}
    \includegraphics[width=0.85\textwidth]{img/tema do tcc.png}
    \fonte{Adaptado de \citeonline{Moretti_ComoEscolherTemaTCC_2024}}
    \label{fig:esquema_metodologico}
\end{figure}

\subsection{POPULAÇÃO E AMOSTRA}
\label{subsec:populacao-amostra}

A população desta pesquisa é constituída por [DEFINIR A POPULAÇÃO]. Para a definição da amostra, será utilizada a técnica de [TIPO DE AMOSTRAGEM], que, segundo \citeonline{hair2019}, é adequada para garantir a representatividade dos dados.

O tamanho da amostra foi calculated considerando [CRITÉRIOS DE CÁLCULO], resultando em [NÚMERO] participantes/elementos.

\subsection{TÉCNICAS DE COLETA DE DADOS}
\label{subsec:coleta-dados}

A coleta de dados será realizada através de [INSTRUMENTOS DE COLETA], que permitirão obter informações sobre [TIPO DE INFORMAÇÕES].

Os instrumentos escolhidos foram:
\begin{itemize}
    \item \textbf{[INSTRUMENTO 1]:} [DESCRIÇÃO E JUSTIFICATIVA];
    \item \textbf{[INSTRUMENTO 2]:} [DESCRIÇÃO E JUSTIFICATIVA].
\end{itemize}

\subsection{ANÁLISE DOS DADOS}
\label{subsec:analise-dados}

Os dados coletados serão analisados utilizando [TÉCNICAS DE ANÁLISE]. Para dados quantitativos, utilizar-se-á estatística descritiva e inferencial. Os dados qualitativos serão submetidos à análise de conteúdo.

\subsection{ASPECTOS ÉTICOS}
\label{subsec:aspectos-eticos}

Esta pesquisa será desenvolvida respeitando os princípios éticos estabelecidos pela Resolução nº 466/2012 do Conselho Nacional de Saúde. [SE APLICÁVEL: O projeto será submetido ao Comitê de Ética em Pesquisa (CEP) antes da coleta de dados]. Todos os participantes assinarão o Termo de Consentimento Livre e Esclarecido (TCLE).


\section{RESULTADOS ESPERADOS}
\label{sec:resultados-esperados}

Os resultados esperados descrevem os produtos, impactos e transformações que a pesquisa pretende alcançar, devendo alinhar-se aos objetivos propostos:

\begin{itemize}
    \item \textbf{Impactos científicos}: [CONTRIBUIÇÃO PARA O AVANÇO TEÓRICO DA ÁREA];
    \item \textbf{Impactos tecnológicos/sociais}: [MUDANÇAS ESPERADAS NO CONTEXTO PRÁTICO];
    \item \textbf{Produtos e entregáveis}: [LISTAR RELATÓRIOS, SOFTWARES OU ARTIGOS GERADOS];
    \item \textbf{Indicadores de acompanhamento}: [DEFINIR MÉTRICAS E METAS DE AVALIAÇÃO].
\end{itemize}


\section{RECURSOS}
\label{sec:recursos}

Esta seção discrimina os recursos necessários para viabilizar a execução do projeto, conforme orientação da NBR 15287.

\subsection{RECURSOS HUMANOS}
\label{subsec:recursos-humanos}

[DESCREVER A EQUIPE ENVOLVIDA: PESQUISADOR PRINCIPAL, ORIENTADOR, COORIENTADOR E COLABORADORES, DESTACANDO AS ATIVIDADES DE CADA UM].

\subsection{RECURSOS MATERIAIS E INFRAESTRUTURAIS}
\label{subsec:recursos-materiais}

[DETALHAR INFRAESTRUTURA DE LABORATÓRIOS, EQUIPAMENTOS, SOFTWARES E INSUMOS NECESSÁRIOS].

\subsubsection{Orçamento}
\label{subsubsec:orcamento}

A \autoref{tab:orcamento} sintetiza os custos previstos por categoria de despesa para a realização da pesquisa.

\begin{table}[htbp]
    \centering
    \caption{Orçamento estimado da pesquisa}
    \label{tab:orcamento}
    \begin{tabularx}{0.85\textwidth}{Xcr}
        \toprule
        \textbf{Item / Descrição} & \textbf{Qtd.} & \textbf{Valor Total (R\$)} \\
        \midrule
        Material de escritório e insumos & 1 & 150,00 \\
        Impressões e cópias de questionários & 1 & 100,00 \\
        Deslocamento e transporte para coleta & 1 & 200,00 \\
        Licenças de software e ferramentas & 1 & 300,00 \\
        \midrule
        \textbf{TOTAL} & & \textbf{750,00} \\
        \bottomrule
    \end{tabularx}
    \fonte{Elaborado pelo autor.}
\end{table}


\section{CRONOGRAMA DE EXECUÇÃO}
\label{sec:cronograma}

O cronograma apresenta a distribuição temporal das atividades previstas para a realização da pesquisa. O \autoref{qua:cronograma} detalha as etapas e seus respectivos meses de execução.

\begin{quadro}[htbp]
    \centering
    \caption{Cronograma de execução da pesquisa}
    \label{qua:cronograma}
    \begin{tabularx}{\textwidth}{|X|c|c|c|c|c|c|}
        \hline
        \textbf{Atividades} & \textbf{Jan} & \textbf{Fev} & \textbf{Mar} & \textbf{Abr} & \textbf{Mai} & \textbf{Jun} \\
        \hline
        Revisão bibliográfica e estado da arte & \checkmark & \checkmark & & & & \\
        \hline
        Elaboração e validação dos instrumentos & & \checkmark & \checkmark & & & \\
        \hline
        Coleta de dados em campo & & & \checkmark & \checkmark & & \\
        \hline
        Processamento e análise dos dados & & & & \checkmark & \checkmark & \\
        \hline
        Redação do relatório final e artigo & & & & & \checkmark & \checkmark \\
        \hline
        Revisão e submissão/defesa & & & & & & \checkmark \\
        \hline
    \end{tabularx}
    \fonte{Elaborado pelo autor.}
\end{quadro}


% ========================================================================
% ELEMENTOS PÓS-TEXTUAIS
% ========================================================================
\input{estrutura/pos_textuais}

% ========================================================================
% FIM DO DOCUMENTO
% ========================================================================
\end{document}